\section{Experiments}

We evaluate six classification methods on Caltech-101~\citep{fei2004learning}: ResNet-50~\citep{he2016deep}, EfficientNet-B2~\citep{tan2019efficientnet}, ViT-B/16~\citep{dosovitskiy2020image}, ConvNeXt-Tiny~\citep{liu2022convnet}, EVA-02-Small~\citep{fang2024eva}, and an SVM with frozen ResNet-18 features~\citep{donahue2014decaf}.

\subsection{Experimental Setup}

\textbf{Dataset.} Caltech-101 contains 9{,}144 images across 102 categories (101 object classes plus background). We apply a stratified 70/15/15 split, yielding 6{,}352 / 1{,}326 / 1{,}466 samples for train / validation / test.

\textbf{Implementation.} All deep models use ImageNet-pretrained weights (EVA-02-Small uses MIM-pretrained weights on IN-22k, fine-tuned on IN-1k) and are fine-tuned with AdamW~\citep{loshchilov2017decoupled} ($\lambda{=}10^{-4}$) using differential learning rates (backbone $\eta_{\text{bb}}{=}10^{-4}$, head $\eta_{\text{head}}{=}10^{-3}$; ViT and EVA-02 use $2{\times}10^{-5}$ / $2{\times}10^{-4}$). We apply linear warmup (5 epochs) followed by cosine annealing over 30 epochs, with early stopping (patience 5). Label smoothing ($\epsilon{=}0.1$), gradient clipping (max norm 1.0), and standard augmentation (random crop, horizontal flip, color jitter, random erasing) are used throughout. The SVM baseline extracts 512-dim features from a frozen ResNet-18 and trains an RBF-kernel SVM ($C{=}1$, $\gamma{=}\text{scale}$). All experiments use batch size 32 on a single NVIDIA GPU.

\textbf{Evaluation metrics.} We report: (1)~overall accuracy, (2)~top-5 accuracy, (3)~macro- and weighted-averaged precision, recall, and F1-score, (4)~per-class accuracy, and (5)~confusion matrices.

\subsection{Overall Accuracy and Top-5 Accuracy}

Tab.~\ref{tab:main_results} summarizes the test-set performance. EfficientNet-B2 achieves the highest overall accuracy (97.20\%) among the evaluated deep models, followed by ResNet-50 and ConvNeXt-Tiny (both 96.93\%), and ViT-B/16 (96.86\%). These four models perform within a narrow 0.34\% range, indicating that with proper fine-tuning and sufficient pretraining, architecture choice has limited impact on this dataset. The SVM baseline reaches 90.18\%, trailing the deep models by approximately 7 percentage points, highlighting the advantage of end-to-end fine-tuning over fixed-feature classification.

Top-5 accuracy exceeds 98\% for all models, with ViT-B/16 and ConvNeXt-Tiny reaching 99.80\%. This suggests that even when the top-1 prediction is incorrect, the true class almost always appears among the top candidates.

\begin{table}[t]
\centering
\caption{\hl{Test-set accuracy on Caltech-101.} Best values are highlighted.}
\label{tab:main_results}
\setlength{\tabcolsep}{6pt}
\small
\begin{tabular}{lcc}
\toprule
Model & Accuracy (\%) & Top-5 (\%) \\
\midrule
ResNet-50          & 96.93 & 99.52 \\
EfficientNet-B2    & \cellcolor{blue!8}97.20 & 99.52 \\
ViT-B/16           & 96.86 & \cellcolor{blue!8}99.80 \\
ConvNeXt-Tiny      & 96.93 & \cellcolor{blue!8}99.80 \\
EVA-02-Small       & --    & -- \\
SVM (ResNet-18)    & 90.18 & 98.09 \\
\bottomrule
\end{tabular}
\end{table}

\subsection{Training Convergence}

Fig.~\ref{fig:acc_loss_epochs}(a) shows validation accuracy as a function of training epoch for all six models. The SVM baseline is shown as a horizontal dashed line since it has no iterative training. Among deep models, ResNet-50 and EfficientNet-B2 converge rapidly, exceeding 95\% by epoch 2--3. ViT-B/16, trained with a lower learning rate ($2{\times}10^{-5}$), exhibits a slower ramp-up but ultimately reaches a competitive accuracy. ConvNeXt-Tiny shows the steadiest improvement curve. Early stopping is triggered at different epochs for each model, reflecting their varying convergence speeds.

Fig.~\ref{fig:acc_loss_epochs}(b) shows the corresponding training loss (solid) and validation loss (dashed) curves. All deep models exhibit rapid loss reduction in the first 2--3 epochs, followed by a gradual decrease. The gap between training and validation loss remains small throughout, indicating that label smoothing and augmentation effectively control overfitting.

\begin{figure*}[t]
\centering
\includegraphics[width=\textwidth]{figures/fig_acc_loss_epochs.pdf}
\caption{\hl{Training convergence.} (a)~Validation accuracy vs.\ epoch. All deep models converge within 10--15 epochs; the SVM (dashed) has no epoch-based training. (b)~Training loss (solid) and validation loss (dashed) vs.\ epoch.}
\label{fig:acc_loss_epochs}
\end{figure*}

\subsection{Precision, Recall, and F1-Score}

Tab.~\ref{tab:detailed_metrics} reports macro- and weighted-averaged precision, recall, and F1-score. For all deep models, weighted scores are slightly higher than macro scores, reflecting strong performance on frequent classes. The gap is most pronounced for SVM: weighted recall (90.18\%) vs.\ macro recall (84.11\%), a 6.07 pp difference, indicating that the SVM struggles disproportionately on minority classes. Among deep models, EfficientNet-B2 achieves the highest macro recall (96.26\%) and macro F1 (96.04\%), suggesting relatively uniform performance across both frequent and rare categories. ConvNeXt-Tiny leads in macro precision (96.44\%).

\begin{table}[t]
\centering
\caption{\hl{Macro and weighted precision, recall, and F1-score} on the test set. Best values per averaging scheme are highlighted.}
\label{tab:detailed_metrics}
\setlength{\tabcolsep}{3.5pt}
\small
\begin{tabular}{lcccccc}
\toprule
 & \multicolumn{3}{c}{Macro} & \multicolumn{3}{c}{Weighted} \\
\cmidrule(lr){2-4} \cmidrule(lr){5-7}
Model & Prec. & Rec. & F1 & Prec. & Rec. & F1 \\
\midrule
ResNet-50       & .9633 & .9563 & .9577 & .9707 & .9693 & .9688 \\
EfficientNet-B2 & .9629 & \cellcolor{blue!8}.9626 & \cellcolor{blue!8}.9604 & \cellcolor{blue!8}.9744 & \cellcolor{blue!8}.9720 & \cellcolor{blue!8}.9718 \\
ViT-B/16        & .9591 & .9555 & .9541 & .9721 & .9686 & .9684 \\
ConvNeXt-Tiny   & \cellcolor{blue!8}.9644 & .9571 & .9586 & .9710 & .9693 & .9689 \\
EVA-02-Small    & --    & --    & --    & --    & --    & -- \\
SVM             & .9389 & .8411 & .8704 & .9314 & .9018 & .9024 \\
\bottomrule
\end{tabular}
\end{table}

\subsection{Per-Class Accuracy}

Fig.~\ref{fig:perclass_all} shows the per-class accuracy for each of the six models, with all 102 classes sorted by their average accuracy across models (ascending). The overall accuracy for each model is indicated by a vertical dashed line. The five deep models achieve near-perfect accuracy on the majority of classes, while the SVM exhibits a substantially longer lower tail, with multiple classes falling below 60\%.

\begin{figure*}[t]
\centering
\includegraphics[width=\textwidth]{figures/fig_perclass_all.pdf}
\caption{\hl{Per-class accuracy for all six models} (102 classes, sorted by average accuracy ascending). Vertical dashed lines mark overall accuracy. Deep models cluster near 100\% on most classes; the SVM shows greater variance.}
\label{fig:perclass_all}
\end{figure*}

\subsection{Confusion Matrices}

Fig.~\ref{fig:cm_all} presents the row-normalized confusion matrices for all six models side by side. For the five deep models, the strong diagonal confirms high per-class recall across the majority of categories; off-diagonal mass is sparse and scattered among visually similar classes. The SVM confusion matrix shows a weaker diagonal with more diffuse off-diagonal entries, particularly for minority classes where the fixed features fail to provide sufficient discriminative power.

\begin{figure*}[t]
\centering
\includegraphics[width=\textwidth]{figures/fig_cm_all.pdf}
\caption{\hl{Normalized confusion matrices} for all six models (102 classes). Each row is normalized to sum to 1. Deep models show strong diagonals; the SVM exhibits more scattered misclassifications.}
\label{fig:cm_all}
\end{figure*}
